\clearpage
\SourceLeaf{362}%
\pdfbookmark[0]{Anzeigen. Mollweide: Archimedes. 1808.}{anzeigen-mollweide-1808}
\begin{center}
\rule{.85\textwidth}{.3pt}\par
Göttingische gelehrte Anzeigen.\quad 1808 Januar 9.\par
\rule{.85\textwidth}{.3pt}
\end{center}
\hspace*{\parindent}\textsc{Archimedes} gründete bekanntlich in seiner Schrift, \textit{Circuli dimensio}, seine Bestimmung der Grenzen für den Umfang des Kreises darauf, dass er denselben zwischen den Umfang eines umgeschriebenen und eines eingeschriebenen 96 Ecks einschloss. Die Berechnung dieser Zahlen, oder vielmehr die Bestimmung einer grössern Zahl, als jener, und einer kleinern, als dieser, verrichtet er durch stufenweises Fortschreiten vom Sechseck zum Zwölfeck, von diesem zum 24 Eck u.\,s.\,f. Für beide 96 Ecke geht er daher, nach unserer Art zu reden, von einem genäherten Werthe der Irrationalgrösse \(\BareRoot{3}\) aus, wovon der eine, nemlich \(\frac{265}{153}\), etwas zu klein, der andere, \(\frac{1351}{780}\). etwas zu gross ist; jener wird bei den umschriebenen, dieser bei den eingeschriebenen Vielecken gebraucht. Bei genauerer Ansicht findet man, dass diese genäherten Werthe in der Reihe \(\frac21,\frac53,\frac74,\frac{19}{11},\frac{26}{15}\) u.\,s.\,f., deren Glieder abwechselnd grösser und kleiner sind als \(\BareRoot{3}\), und jedes weniger davon verschieden, als irgend ein andrer, durch kleinere Zahlen ausgedrückter, Bruch, --- mit vorkommen; der Bruch \(\frac{265}{153}\) ist nemlich das achte, und \(\frac{1351}{780}\) das eilfte Glied der Reihe. Es scheint demnach, dass \textsc{Archimed} diese genäherten Werthe nicht durch Zufall, sondern methodisch gefunden habe; da er selbst sich aber über die Art, wie er dazu gekommen ist, gar nicht erklärt, und man übrigens nicht findet, dass unsre Methoden dergleichen Aufgaben aufzulösen, den Alten bekannt gewesen wären, so bietet sich hier ein Gegenstand zu Conjecturen dar. Hr. Prof. \textsc{Mollweide} in Halle hat in einer kürzlich an die Königl. Societät, deren Correspondent er ist, eingeschickten kleinen Abhandlung, welche
\begin{center}\textit{De methodo ab Archimede adhibita ad rationem, in qua inter se sunt latus trianguli aequilateri et radius circuli circumscripti, numeris veritati proxime exprimendam}\end{center}
\noindent überschrieben ist, eine Untersuchung angestellt, und ein Verfahren angegeben, das dem Zustande der Arithmetik der Alten angemessen ist, und also vielleicht das von \textsc{Archimed} gebrauchte selbst sein könnte. Hr. M. leitet nemlich, indem er die Seite des Dreiecks durch \(AC\), und den Halbmesser des umschriebenen
